\scanpage{70}
$E^{2n}(S)=$ ens des classes de couples $(E,F)$, $E$ et $F$ formes quadratiques \unsure{admissibles} de \unsure{rangs} \deleted{\unsure{2n}} \inserted{\uuline{pairs}} telles que $\operatorname{rang}E-\operatorname{rang}F=$
\[
x_1^2+\text{---}+x_n^2+\deleted{xy}\;y_1y_2+y_3y_4\text{---}+y_{2m-1}y_{2m}
\]
$\bar E^{n,m}=$ ens des classes à isom près des formes quadratiques \unsure{pointées} de type $(n,m)$
\[
\coprod_{n,m\geq0}\bar E^{n,m}(S),=\bar E(S)
\]
est muni d’une structure de \uline{monoïde} \uline{bigradué}.

\deleted{Soit} [N.B. Si $2$ est \unsure{inv.} alors $\bar E^{\unsure{n},m}\simeq\bar E^{\deleted{\unsure{m}}n+2m,0}$

Si $2=0$, \deleted{$\unsure{A\text{---}0}$}, alors $\bar E^{n,m}\simeq\bar E^{1,m}$ ;

\deleted{\unsure{mais}}

en particulier, \unsure{on ne peut retrouver} : \deleted{\unsure{l’isom.}}
\[
\unsure{E\simeq}E^{\alpha,0}\quad\text{ou}\quad E^{1,\beta},\quad\text{ou}\quad E^{0,\unsure{\gamma}}\text{---}\quad]
\]
$E(S)$ groupe associé au monoïde
\[
\deleted{\coprod_{n,m\in\mathbf Z}\bar E^{n,m}(S)}
\]
\diagram{\begin{tikzcd}
E(S)\arrow[r,"\epsilon_1\times\epsilon_2"]&\mathbf Z\times\mathbf Z
\end{tikzcd}}
\[
E^{n,}
\]
\margin{69}
